% A.123: complete harmonic octic evaluation and finite initialization.
\subsection{Full harmonic initialization at exponent eight}

\begin{statement}{A.123}{Proved}
{Every common harmonic octic character has zero normalized local evaluation}
Put $\mathcal C=\{0,1,\tfrac12\}$.
Let $R\in\Q\setminus\mathcal C$ and consider the actual
pointed eighth-power cover
\[
 C_R^{(8)}:\quad
       T-a=(R-a)X_a^8\quad(a\in\mathcal C).
\tag{A.123.1}
\]
For $c\in\Q\setminus(\mathcal C\cup\{R\})$, let $B_8$
be the complete common abelian field of
\[
 \mathcal E_{8,c}=
  \Q[Z_a^{\pm1}:a\in\mathcal C]
       \big/\left(Z_a^8-\frac{c-a}{R-a}\right)_{a\in\mathcal C},
\tag{A.123.2}
\]
computed in A.121.
For every character $\chi$ of
$\operatorname{Gal}(B_8/\Q)$, every place $v$ of $\Q$,
and every ordinary point of $C_R^{(8)}(\Q_v)$ with
$T\ne c$, the actual cyclic-algebra evaluation satisfies
\[
                    \left(\frac{T-c}{R-c},\chi\right)=0
                                      \quad\text{in }\operatorname{Br}(\Q_v).
\tag{A.123.3}
\]
This assertion includes every positive and negative
mixed fourth-power relation in A.121.

Consequently, for any rational new puncture
$c\notin\mathcal C$, any finite set $S$ of places, and
one prescribed ordinary point of (A.123.1) at each
$v\in S$ avoiding $T=c$, there is a rational coefficient
$a_c\ne0$ for which
\[
 \begin{gathered}
 T-a=(R-a)X_a^8\quad(a\in\mathcal C),\qquad
 T-c=a_cY^8,\\
 (T,(X_a)_a,Y)\longmapsto(T,(X_a)_a)
 \end{gathered}
\tag{A.123.4}
\]
has ordinary points over every completion and lifts
each prescribed old point.
This is the full finite exponent-eight initialization
on the harmonic four-punctured source of A.107.
It does not prescribe all infinitely many old local
points or assert compatibility between coefficients
chosen at different exponents.
\end{statement}
\begin{proof}
Write $D_a=(c-a)/(R-a)$ and $f=(T-c)/(R-c)$.
The case $T=R$ gives $f=1$, so we may exclude it.

\emph{The local field-factor and norm maps.}
For $k=\Q_v$, restrict the character $\chi$ to $G_k$,
and let $K_\chi/k$ be its actual cyclic field.
It has degree dividing four by A.121.
Every field factor of $\mathcal E_{8,c}\otimes k$
contains $K_\chi$.
Indeed an embedding of such a factor into $\overline k$
selects an actual root tuple.
Its local stabilizer is contained in the global
root-tuple stabilizer on which $\chi$ vanishes.
This proves the claimed field inclusion, including
every local factor and every root choice.

We will repeatedly use the following literal norm map.
If all the $D_b$ except one, say $D_a$, are eighth
powers in $k$, choose their actual roots $u_b\in k$.
Specializing $Z_b\mapsto u_b$ gives a surjection
\[
 \mathcal E_{8,c}\otimes k
             \longrightarrow k[Z_a]/(Z_a^8-D_a).
\tag{A.123.5}
\]
Thus $K_\chi$ lies in every field factor of the target.
The polynomial norm identities, with one root chosen
in each factor, give
\[
 X^8-D_a=\prod_E N_{E/k}(X-\theta_E),\qquad
 1-D_a=\prod_E N_{E/k}(1-\theta_E).
\tag{A.123.6}
\]
Each norm on the right factors through $K_\chi$.
Whenever $X^8=(T-a)/(R-a)$, their ratio is exactly $f$.
The numerator and denominator are nonzero because
$T,c,R$ are distinct from the relevant old puncture
and $c\ne R,T$.
This proves (A.123.3) in the situation of (A.123.5).

\emph{The exact harmonic coordinate maps.}
The maps $g_1(z)=1-z$ and $g_2(z)=1/(2z)$ permute
the four punctures.
They act on the old root coordinates, in the order
$0,1,\tfrac12$, by
\[
 \begin{array}{c|ccc}
        &Y_0&Y_1&Y_{1/2}\\ \hline
  g_1&X_1&X_0&X_{1/2}\\
  g_2&X_0^{-1}&X_{1/2}/X_0&X_1/X_0.
 \end{array}
\tag{A.123.7}
\]
The same formulas with $Z$ in place of $X$ are
isomorphisms of the actual boundary-fibre algebras.
For $g_2$ the transformed radicands and normalized
function are
\[
 (D'_0,D'_1,D'_{1/2})
       =(D_0^{-1},D_{1/2}/D_0,D_1/D_0),
                 \qquad f'=fR/T.
\tag{A.123.8}
\]
For $g_1$ the first two radicands are exchanged
and $f'=f$.
The inverse root maps are obtained by the same
involutions.
Hence the common abelian field and its actual
characters are unchanged.
The factor $R/T=X_0^{-8}$ in (A.123.8) is an eighth
power, which is a norm from every cyclic field of
degree dividing four.
Therefore these maps preserve the cyclic-algebra
evaluation for the same character $\chi$.

At a single local place, we may also exchange $R,T$.
Dividing every boundary root by its actual old root
$X_a$ identifies the new radicands
$D'_a=D_a/X_a^8$ with the same local fibre algebra.
The normalized entry becomes $f^{-1}$.
Thus the exchange preserves vanishing, with the
local character retained.
This last operation is used only in the odd-place
argument below.

\emph{The two quartic reductions at the dyadic place.}
Now set $k=\Q_2$.
Let $K=\Q(\sqrt{D_a}:a\in\mathcal C)$ and let $B_4$
be the corresponding quartic common field.
The actual map $C_R^{(8)}\to C_R^{(4)}$ is
$X_a\mapsto X_a^2$.
By A.112, every character in $B_4$ has zero normalized
evaluation at this point.
If $i\notin K$, A.121 gives $B_8=K=B_4$; if
$i\in K$ but $\sqrt2\notin B_4$, it gives $B_8=B_4$.
Both cases are therefore finished.
We may assume
\[
                  \mu_8\subset B_4\subset B_8.
\tag{A.123.9}
\]
In particular $i\in K$ and the two actual evaluations
$(f,-1)$ and $(f,2)$ are zero.

Apply the rational coordinate maps (A.123.7) until
$e=v_2(R)\ge1$.
For clarity, if $v_2(R)=0$, apply $g_1$; if
$v_2(R)<0$, apply $g_2$, whose new valuation is
$-1-v_2(R)\ge0$, and apply $g_1$ once more when this
is zero.
The excluded values of $R$ ensure none of these
maps reaches an old puncture.
Write $e_T=v_2(T)$ and $s=v_2(T-R)$.
The old eighth-power ratios give
\[
 \begin{gathered}
 v_2(T)\equiv e,\quad v_2(T-1)\equiv0,\quad
                   v_2(T-\tfrac12)\equiv-1\pmod8,\\
 e_T>0,\qquad s\ge5,\qquad e_T\equiv e\pmod8.
 \end{gathered}
\tag{A.123.10}
\]
Here negative $e_T$ gives $v_2(T-1)=e_T$, so the
second congruence first forces $8\mid e_T$.
Then $e_T\le-8$ and $v_2(T-\tfrac12)=e_T$,
contradicting the third congruence.
If $e_T=0$, both $e$ and $v_2(T-1)$ would be positive
multiples of eight.
Then $(1-2T)/(1-2R)\equiv-1\pmod{32}$, contrary
to its being an eighth-power unit.
Finally $(1-T)/(1-R)\in1+32\Z_2$ gives $s\ge5$.

\emph{Equal old valuation coordinates.}
If $e_T=e$, all three old valuation coordinates agree:
they are $(e,0,-1)$.
More generally, for any finite set of old punctures
with equal valuation coordinates at $R,T$, put
$M=\max_a v_2(R-a)$.
The unit eighth-power ratios give
\[
                           s\ge M+5.
\tag{A.123.11}
\]
Since $i\in K$, there is a subset of the $D_a$
whose product is a negative rational square.
If $v_2(c-R)\ge M+2$, every
\[
                 D_a=1+\frac{c-R}{R-a}
\]
lies in $1+4\Z_2$.
Their subset product then cannot be a negative
unit square, which is $3$ modulo four.
We have proved $v_2(c-R)\le M+1$, and hence
\[
                          f\in1+16\Z_2.
\tag{A.123.12}
\]
Every such dyadic unit is an actual fourth power.
For example the square map sends $1+8\Z_2$ onto
$1+16\Z_2$ and sends $1+4\Z_2$ onto $1+8\Z_2$,
by Hensel lifting the respective normalized
polynomials.
Thus (A.123.12) is a norm from every cyclic field
of degree dividing four and proves the desired
evaluation.
This deals in particular with every shallow case
$1\le e\le4$: an unequal depth would make
$s=\min(e,e_T)<5$, contrary to (A.123.10).

\emph{Different depths in the same dyadic boundary disc.}
We are left with $e_T\ne e$.
Equation (A.123.10) now gives the exact bounds
\[
                 e,e_T\ge5,\qquad
                 s=\min(e,e_T),\qquad e\equiv e_T\pmod8.
\tag{A.123.13}
\]
Put $k=v_2(c)$.
If $k\le0$, then $v_2(R-c)=k$ and
$v_2(f-1)=s-k\ge5$, so $f$ is a fourth power.

Suppose next that $1\le k\le4$.
Again $v_2(R-c)=k$; put $d=s-k\ge1$.
If $d\ge4$, $f$ is a fourth power.
The two known quartic-character norm conditions in
(A.123.9) exclude $d=1,2$.
Here is their exact local residue test.
A unit norm from $\Q_2(i)$ is $1$ modulo four:
in $a^2+b^2$, valuation zero forces $a,b$ integral
with exactly one a unit.
A unit norm from $\Q_2(\sqrt2)$ is $1$ or $7$
modulo eight: in $a^2-2b^2$, the valuations of the
two terms have different parity, so valuation zero
forces $a$ a unit and $b$ integral, giving those
two residues.
Consequently a unit that is both norms is $1$
modulo eight.
For $d=1$, $f$ is $3$ modulo four; for $d=2$,
it is $5$ modulo eight.
Both possibilities contradict the proved norms.

It remains in this range that $d=3$.
Now $f\in1+8\Z_2$ is a square.
We prove that every common character is quadratic
in precisely this remaining configuration.
Since $s\ge5$, one has $k=s-3\ge2$.
Moreover
\[
 \begin{gathered}
 v_2(D_0)=k-e\equiv-3\pmod8,\\
 D_1=\frac{1-c}{1-R}\in1+4\Z_2,\qquad
 D_{1/2}=\frac{1-2c}{1-2R}\in1+4\Z_2.
 \end{gathered}
\tag{A.123.14}
\]
Take any relation $t\in\mathcal R_4$ in the exact
list A.121.
Its product has valuation divisible by four.
Because the first valuation in (A.123.14) is odd,
and $t_0\in\{0,1,2,3\}$, this forces $t_0=0$.
The product then lies in $1+4\Z_2$ and cannot equal
$-r_t^4$, whose unit residue is $-1$ modulo four.
Thus every such relation has $\delta_t=0$.
A.121 now expresses $B_8$ using only the quadratic
field $B_4$ and the positive-relation quadratic
fields $\Q(\sqrt{r_t})$.
Every character is quadratic and the square $f$
kills its evaluation.
This retains all mixed relation vectors; no
restriction to subset relations has been made.

Finally suppose $e,k\ge5$.
Both
\[
 D_1=\frac{1-c}{1-R},\qquad
 D_{1/2}=\frac{1-2c}{1-2R}
\tag{A.123.15}
\]
lie in $1+32\Z_2$, and therefore have actual eighth
roots in $\Q_2$.
The specialization (A.123.5) leaves the single
factor polynomial $Z_0^8-D_0$.
Equations (A.123.6), with $X=X_0$ and
$D_0=c/R$, express
\[
                         f=\frac{X_0^8-D_0}{1-D_0}
\tag{A.123.16}
\]
as a norm from $K_\chi$ for every common character.
This proves every dyadic case.

\emph{Every odd finite place.}
Let $k=\Q_p$ with $p$ odd.
The four punctures have distinct reductions.
Call a point good if its reduction avoids them,
and otherwise specify its boundary residue disc.
The transformations (A.123.7) lie in
$\operatorname{PGL}_2(\Z_p)$, preserve good points,
and act transitively on the four boundary discs.
Every principal unit is an eighth power by Hensel,
since eight is invertible.

If $R,c$ lie in the same boundary disc, send it to
zero.
Then $D_1,D_{1/2}$ are principal units, so
(A.123.5)--(A.123.6) kill every common character.
If $R,T$ lie in the same boundary disc and $c$
lies elsewhere, send that disc to zero.
The identity
\[
                  f=\frac{1-T/c}{1-R/c}\in1+p\Z_p
\tag{A.123.17}
\]
makes $f$ an eighth power.

If exactly one of $R,T$ is good, exchange them
by the local fibre map following (A.123.8), so
that $R$ is good.
Suppose $R$ is good.
If $c$ is in a boundary disc, send it to zero.
If $T$ is also in that disc, exchange the two old
points and apply the first case.
Otherwise $c$ has positive valuation, $R$ is a
unit, and $T$ has nonpositive valuation.
Then
\[
                    f/(T/R)=\frac{1-c/T}{1-c/R}
                                      \in1+p\Z_p,
\tag{A.123.18}
\]
so $f$ is an eighth power.

Suppose instead that $R,c$ are both good.
All three $D_a$ are units.
Every root of every $Z^8-D_a$ lies in an
unramified extension: take the residue-field
roots in a finite extension and lift them by
Hensel.
Thus the actual local splitting algebra and
each $K_\chi$ are unramified.
If $R$ reduces to $c$, all $D_a$ are principal
units and the fibre has a $k$-rational root tuple.
If $T$ is good and reduces to $c$, each $D_a$
is the old eighth-power ratio multiplied by a
principal unit, giving the same conclusion.
A common character must then be locally trivial.
In all other configurations $v_p(f)=0$, except
when $T$ is in the disc of infinity.
There $(T-1)/(R-1)$ gives $v_p(T)\in8\Z$ and
$v_p(f)=v_p(T)$.
Every unit, and every element of valuation
divisible by eight, is a norm from an unramified
cyclic extension of degree dividing four.
The residue-field norm is onto, its principal-unit
norm is onto because the base has the required
degree-th roots of all principal units, and the
norm of the base uniformizer is its degree-th
power.
This proves the evaluation in this case.

The only remaining configuration has $R,T$ in
different boundary discs.
If $c$ shares either disc, the first case and
an exchange apply.
If $c$ is in a third boundary disc, send it to
zero and use (A.123.18).
If $c$ is good, send $R$ to zero and put
$e=v_p(R)>0$.
For $T$ in another finite boundary disc,
$T/R\in k^{\times8}$ gives $e\in8\Z$.
For $T$ in the disc of infinity,
$(T-1)/(R-1)$ first gives $v_p(T)\in8\Z$,
and then $T/R$ again gives $e\in8\Z$.
Here $D_0$ has valuation $-e\in8\Z$ and the
other $D_a$ are units.
After scaling the first root by $p^{e/8}$,
all its roots also lie in an unramified extension.
The character $K_\chi$ is therefore unramified,
while $v_p(f)$ is zero or $v_p(T)\in8\Z$.
The same norm criterion applies.
These cases exhaust all odd-place configurations.

\emph{The real place and the finite initialization.}
The three positive eighth-power ratios put $R,T$
in the same connected interval of
$\R\setminus\{0,\tfrac12,1\}$.
If $c$ does not lie strictly between them, then
$f>0$, a norm from $\C$.
If $c$ lies strictly between them, each $D_a>0$
and the fibre has a real root tuple.
Every common character is then zero on $G_\R$.
This proves (A.123.3) at infinity as well.

For $c\ne R$, subtracting the constant global
class $(R-c,\chi)$ makes every local evaluation
in the finite-character criterion A.85 zero by
(A.123.3).
The constant's invariant sum is zero by global
reciprocity.
That exact criterion supplies the rational
coefficient and the forgetful morphism (A.123.4),
including the prescribed finite collection of
old points.
For $c=R$, the fibre has the rational root tuple
$(1,1,1)$ and therefore has common abelian field
$\Q$; the same criterion has no nonzero character
to test.
This proves the finite exponent-eight assertion.
\end{proof}

The certificate
\path{code/harmonic_eighth_full_certificate.py}
checks the literal fibre-algebra changes, the dyadic
norm residue restrictions, all finite depth and mixed
relation implications in the proof, and exact local
evaluations of every character supplied by A.121 on
the retained finite collection of harmonic fibres.
Its output is retained beside the program.
The general all-point assertion follows from the
preceding proof, with all field and norm maps displayed.

